%! Constructing a Confidence Interval for a Population Proportion — Unit 6, Lesson 6.2 — Student Packet Cover Sheet
\documentclass[10pt]{article}
\usepackage{apstats-article}
\usepackage{apstats-boxes}
\usepackage{ltablex}
\keepXColumns

\begin{document}

% Full-bleed navy banner
\begin{tikzpicture}[remember picture, overlay]
  \fill[navy] (current page.north west)
    rectangle ([yshift=-0.9in]current page.north east);
\end{tikzpicture}
\vspace{-0.8in}

\begin{center}
  {\color{white}\textbf{\LARGE AP Statistics}}\\[4pt]
  {\color{skymid}\large\textbf{Unit 6}}\\[6pt]
  {\color{white}\textbf{\large Lesson 6.2 \quad Constructing a Confidence Interval for a Population Proportion}}
\end{center}

\vspace{0.22in}
\namedateperiod
\vspace{0.18in}

\begin{learningtargetbox}
\small
\begin{enumerate}[leftmargin=1.5em, itemsep=5pt]
  \item I can verify the Random, Large Counts, and 10\% conditions for a one-sample $z$-interval for a proportion. \textit{(UNC-4.B, Skill 4.C)}
  \item I can construct and interpret a one-sample $z$-interval for a population proportion. \textit{(UNC-4.A/D/E, Skill 3.D)}
  \item I can determine the sample size needed to achieve a specified margin of error. \textit{(UNC-4.C, Skill 3.D)}
\end{enumerate}
\end{learningtargetbox}

\vspace{0.14in}

\begin{tocbox}
\small
\renewcommand{\arraystretch}{1.5}
\begin{tabularx}{\linewidth}{@{} c l X r @{}}
\rowcolor{sky}
\textbf{\#} & \textbf{Component} & \textbf{Description} & \textbf{Score} \\
\midrule
1 & Warm-Up        & Sampling distributions \& z-scores review & \blank{1.2cm} \\
2 & Guided Notes   & One-sample $z$-interval procedure          & \blank{1.2cm} \\
3 & Group Activity & Build and interpret CIs from real scenarios & \blank{1.2cm} \\
4 & Exit Ticket    & Full procedure from single scenario        & \blank{1.2cm} \\
5 & Homework       & Practice: three intervals + sample size    & \blank{1.2cm} \\
\midrule
  & \multicolumn{2}{r}{\textbf{Total}} & \blank{1.2cm} \\
\end{tabularx}
\end{tocbox}

\end{document}
